\section{Conditions on a mechanism}\label{sec:conditions}

A candidate arising from this picture inherits the necessary conditions the
companion investigation has accumulated \cite{Selection}. We do not restate their
derivations; the table records which bear on a flow, and what this paper has
added to each.

\begin{center}
\begin{tabular}{@{}p{0.28\textwidth}p{0.64\textwidth}@{}}
\toprule
Inherited condition & Bearing, after this paper\\
\midrule
Reflection invariance & The flow must commute with $x\mapsto-x$; the involution is
the functional equation, which is also what makes
Proposition~\ref{prop:unimodular} true.\\
Unboundedness relative to $L^2$ & Now a statement about \emph{time}: by
Proposition~\ref{prop:groupdelay} the delay grows like $\log\tau$, so a packet at
frequency $N$ is held for $\sim\log N$. Condition~\ref{cond:bilinear} restates the
exclusion in the form a candidate can fail.\\
Spectral weight and channel tower & The archimedean multiplier \emph{is} the group
delay (Proposition~\ref{prop:groupdelay}), so the density identity is no longer an
input to be matched but a property of a phase.\\
Criticality & Section~\ref{sec:region}: the flow formulation inherits it. A bound
at fixed $\omega$ buys a zero-free strip of width $\omega$ and no more.\\
Jump structure & Realized on the connection by \eqref{eq:atomic}: the atoms of the
kernel are the prime periods.\\
\bottomrule
\end{tabular}
\end{center}

Three conditions are specific to the deformation picture. They are
\emph{proposals}: each is elementary, none has been tested against a specific
theory, and the third carries an identification hypothesis this program
deliberately does not make.

\subsection{The object is a bilinear, not a character}

\begin{condition}[Bilinear endpoints, and a delay test]\label{cond:bilinear}
The object must be a bilinear in unbounded endpoint operators joined by a
transport --- not a character --- and its group delay must grow like $\log\tau$
with smooth part $2\theta'(\tau)$.
\end{condition}

The first clause names a class and the second is the test. The exclusion behind
the clause is worth restating precisely, because it is narrower than it looks. A
holonomy valued in a compact group, in a finite-dimensional representation, is
unitary, so its character is a bounded function and every pairing assembled from
it is bounded on $L^2(I_L)$ --- excluded by the unboundedness rule. That argument
uses only that the object is a \emph{trace of a group element}. It is also the
structural reason the conformal Wilson benchmark of \cite{Selection} failed:
$2\cos\theta$ is bounded \emph{because} it is a character of $SU(2)$ in the
fundamental, not by accident of that benchmark.

An \emph{open} Wilson line is not a character. In the construction of
\cite{OWL} the operator is
\[
 OW_i^{\,j}[\tilde C]=\bar\psi^{\,j}(x_j)\;
 P\exp\!\left[\int_{\tilde C}\!ds\,\Big(iA_\mu\dot x^\mu
 +n^k(s)\Phi_k\,|\dot x|\Big)\right]\psi_i(x_i),
\]
a bilinear in matter fields joined by a Wilson line, the endpoint fields
transforming in the fundamental and antifundamental so that the contraction is
gauge invariant. Its endpoints carry operator-valued distributions, so the object
is unbounded by construction and the character argument has nothing to act on.
\textbf{The condition is satisfied structurally rather than evaded}, which is why
it is stated here as a prescription rather than a prohibition.

Two further remarks connect that construction to Section~\ref{sec:endpoint}.
First, the geometry matches: the supersymmetric configuration of \cite{OWL} in
the defect theory is a \emph{ray ending on a defect}, one end pinned and one end
free, and Proposition~\ref{prop:endpoint} found the variation to be rank one at
the leading end only, the trailing one killed by causality --- with the boundary
value of the evolved field as the coefficient, which in the open Wilson line is
the endpoint field itself. Second, the obstruction of Proposition~\ref{prop:trace}
is, in that language, the ordinary fact that a field at a point is a distribution
and must be smeared; Condition~\ref{cond:marginal} fixes the profile.

Note also that $V_{\omega,L}$ is not unitary --- it is a contraction and the
object of interest is its defect \eqref{eq:defect} --- so the transport to be
matched is not anti-Hermitian either.

\subsection{The non-commutativity must be made visible}

\begin{condition}[Visible transverse structure]\label{cond:visible}
A mechanism must supply a reason for the holonomy to see the non-commutativity of
Proposition~\ref{prop:algebra}: either a domain that is not simply connected, or a
failure of the connection to exist.
\end{condition}

This replaces the weaker demand of Remark~\ref{rem:cost}, which asked only for
non-commutativity transverse to $\omega$. That exists
(Proposition~\ref{prop:algebra}); what it lacks is visibility
(Proposition~\ref{prop:flat}). Problems~\ref{prob:trace} and \ref{prob:curvature}
are the two forms the condition can take.

\subsection{Marginal short distance: effective dimension one half}

\begin{condition}[Marginality]\label{cond:marginal}
Suppose the line's parameter is identified with the arithmetic coordinate $x$.
Since $\gammak(r)=\frac1{2r}+O(1)$ as $r\downarrow0$ by \eqref{eq:gamma-kernel},
\eqref{eq:gamma-energy} gives
\[
 K[F]=\frac12\iint|F(x)-F(y)|^2\gammak(|x-y|)\dd x\dd y
 \sim\frac14\iint\frac{|F(x)-F(y)|^2}{|x-y|}\dd x\dd y
\]
near the diagonal: the Gagliardo seminorm at the borderline exponent $s=0$. An
insertion on the line reproducing this must have two-point function $|x-y|^{-1}$,
an effective defect dimension $\Delta=\frac12$, with the coefficient fixed.
\end{condition}

This excludes the first thing one would try. The displacement operator of a line
defect has protected dimension $\Delta=2$, fixed by the Ward identity that defines
it, and the standard protected scalar insertions on a supersymmetric Wilson line
have $\Delta=1$ \cite{GiombiKomatsu,CHMS}. \emph{So the deformation whose
generator matches is not the displacement operator}, which is what a deformation
of the contour would naively produce.

\begin{remark}[The principal series, and why everything here is marginal]
\label{rem:principal}
There is one reading under which Condition~\ref{cond:marginal} is not an accident.
For a one-dimensional conformal system the dilatation eigenvalue is the conformal
weight, and the unitary principal series of $SL(2,\R)$ is $\Delta=\frac12+i\tau$
with $\tau$ real --- the critical line. In that reading the zeros are
principal-series data, the shift $\omega$ moves off the unitary line, which is the
detection threshold of Corollary~\ref{cor:graded}, and the Riemann hypothesis is
the statement that nothing leaves it. Condition~\ref{cond:marginal}'s
$\Delta=\frac12$ is then the \emph{edge} of the principal series, the most
marginal unitary weight there is, which would explain the coincidence of
Remark~\ref{rem:marginal}: the Gagliardo borderline, the transfer's
$(2+|\tau|)^{-1/2}$ gain and the trace threshold $r>\frac12$ are three faces of
one exponent. \emph{This is a reading, not a result. It rests on the
identification the program declines to make, and the neighbourhood --- conformal
quantum mechanics and the zeros --- is crowded.}
\end{remark}

\subsection{Conformal invariance is a liability, not an asset}

It is natural to ask whether the conformal symmetry of Yang--Mills could be used
to compactify and tame the infinities of the problem. The answer separates.

\emph{What compactifies.} The spectral infinity does, and the compactification is
already in hand: the conformal map of the right half-plane to the disc is the
Cayley transform of Proposition~\ref{prop:cayley}, under which $K_\omega$ becomes
an inner function of the disc whose zeros accumulate at one boundary point, the
logarithmically divergent delay of Proposition~\ref{prop:groupdelay} being exactly
that accumulation. The length infinity also does: the conformal map of a ray to a
semicircle \cite{OWL} turns $L\to\infty$ into the closing-up of an arc.

\emph{What does not.} Proposition~\ref{prop:trace} is a short-distance statement
and is invariant under reparametrization; smearing, not compactification, is its
fix. And Corollary~\ref{cor:graded} is a statement about the zero set of $\xi$,
so compactifying moves the content to a boundary point without removing it ---
which matters here, because by Section~\ref{sec:region} the content is
\emph{entirely} at that boundary point.

\emph{What it costs.} Exact scale invariance is incompatible with the target. A
scale-invariant theory has no absolute scale, and the Weil form fixes one twice
over: dilation destroys the prime atoms at $m\log p$ and shifts the density
symbol. This is why the renormalization-group fixed point was closed in
\cite{ExclusionMap}. The open Wilson line construction makes the tension concrete
in one class: \cite{OWL} finds that in non-superconformal $\mathcal N=2$ theories
with fundamental hypermultiplets \emph{all BPS open Wilson lines have vanishing
expectation value}, nontrivial protected operators appearing only in the
superconformal case. In this program's terms,

\begin{center}
\emph{protection requires conformality, conformality forbids an absolute scale,
and the target requires one.}
\end{center}

That inverts the preference recorded in the program's defect survey, which ranks
supersymmetric cases higher because they offer more explicit identities: here the
explicit identities are available exactly where the required structure cannot be.
Pure Yang--Mills, whose dimensional transmutation manufactures an absolute scale,
has the feature the target needs and none of the identities.

\emph{Where a scale could come from.} Not from a single conformal defect, which
preserves dilations about a point on it. Scales are \emph{separations}: in the
radial coordinate $r=e^x$ of a ray, the prime atoms sit at $r=p^m$, so the
structure required is invariance under $r\mapsto pr$ for every prime and not under
continuous dilation. That is \eqref{eq:atomic} in geometric dress, with the
punctures at integer radii, and it is the same statement as
Condition~\ref{cond:visible}. The named object in that direction is the
Bost--Connes system \cite{BostConnes}; we record where the symmetry lives and do
not enter that literature.
